\chapter{多电子原子}
\problem{5.1}{$\mathrm{He}$ 原子的两个电子处在 $2\mathrm p3\mathrm d$ 电子组态。问可能组成哪几种原子态？用原子态的符号表示之。已知电子间是 $LS$ 耦合。}
\answerarea
\problem{5.2}{已知 $\mathrm{He}$ 原子两电子被分别激发到 $2\mathrm p$ 和 $3\mathrm d$ 轨道，其所构成的原子态为 ${}^3\mathrm D$。问这两电子的轨道角动量 $\boldsymbol p_{l_1}$ 与 $\boldsymbol p_{l_2}$ 之间的夹角，自旋角动量 $\boldsymbol p_{s_1}$ 与 $\boldsymbol p_{s_2}$ 之间夹角分别是多少？}
\answerarea\clearpage
\problem{5.3}{锌原子（$Z=30$）的最外层电子有两个，基态时的组态是 $4\mathrm s4\mathrm s$。当其中有一个被激发，考虑两种情况：
\begin{enumerate}
\item 那电子被激发到 $5\mathrm s$ 态；
\item 它被激发到 $4\mathrm p$ 态。
\end{enumerate}
试求出在 $LS$ 耦合情况下这两种电子组态分别组成的原子状态。画出相应的能级图。从（1）和（2）情况形成的激发态向低能级跃迁分别各有几种光谱跃迁？}
\answerarea
\problem{5.4}{试以两个价电子 $l_1=2$ 和 $l_2=3$ 为例证明，不论是 $LS$ 耦合还是 $jj$ 耦合都给出同样数目的可能状态。}
\answerarea\clearpage
\problem{5.5}{（原书带 $*$）利用 $LS$ 耦合、泡利原理和洪德定则来确定碳 $Z=6$、氮 $Z=7$ 原子的基态。}
\answerarea
\problem{5.6}{已知 $\mathrm{He}$ 原子的一个电子被激发到 $2\mathrm p$ 轨道，而另一个电子还在 $1\mathrm s$ 轨道。试作出能级跃迁图来说明可能出现哪些光谱线的跃迁。}
\answerarea\clearpage
\problem{5.7}{$\mathrm{Ca}$ 原子的能级是单层和三重结构，三重结构中 $J$ 大的能级高。其锐线系的三重线的频率 $\nu_2>\nu_1>\nu_0$，其频率间隔为 $\Delta\nu_1=\nu_1-\nu_0$，$\Delta\nu_2=\nu_2-\nu_1$。试求其频率间隔比值 $\Delta\nu_2/\Delta\nu_1$。}
\answerarea
\problem{5.8}{$\mathrm{Pb}$ 原子基态的两个价电子都在 $6\mathrm p$ 轨道，若其中一个价电子被激发到 $7\mathrm s$ 轨道，而其价电子间相互作用属于 $jj$ 耦合。问此时 $\mathrm{Pb}$ 原子可能有哪些状态？}
\answerarea\clearpage
